\chapter{Laboratorio ed Esercitazioni}

In questo capitolo vengono raccolti gli snippet di codice essenziali (in ambiente MATLAB) relativi alle esercitazioni del corso. Questi esempi formano la base pratica per lo sviluppo del progetto d'esame.

\section{Pipeline Analisi ECG}
L'elaborazione dell'elettrocardiogramma richiede la pulizia del segnale prima dell'individuazione dei picchi R.
\begin{lstlisting}[language=Matlab, frame=single, commentstyle=\color{green!50!black}, keywordstyle=\color{blue}, stringstyle=\color{purple}]
% 1. Rimozione della deriva di linea di base
ecg_detrend = detrend(ecg_raw, 'linear');

% 2. Filtraggio passa-banda
% (es. per rimuovere rumore muscolare)
[b, a] = butter(3, [0.5 40]/(fs/2), 'bandpass');
% filtfilt evita sfasamenti temporali
ecg_filt = filtfilt(b, a, ecg_detrend);

% 3. Individuazione dei picchi R (complesso QRS)
% MinProminence e' essenziale per ignorare
% picchi spuri come le onde T o P
[peaks, locs] = islocalmax(ecg_filt, 'MinProminence', 0.5);

figure; hold on;
plot(t, ecg_filt);
plot(t(locs), ecg_filt(locs), 'r*', 'MarkerSize', 8);
title('ECG con Picchi R');
\end{lstlisting}

\section{Pipeline Analisi EDA / GSR}
Per l'attività elettrodermica si rende necessaria la separazione della componente tonica (SCL) da quella fasica (SCR).
\begin{lstlisting}[language=Matlab, frame=single, commentstyle=\color{green!50!black}, keywordstyle=\color{blue}]
% 1. Estrazione Componente Tonica (SCL)
% - Filtro Passa-Basso molto stretto
fc_scl = 0.05; % Hz
[b_ton, a_ton] = butter(2, fc_scl/(fs/2), 'low');
GSR_tonica = filtfilt(b_ton, a_ton, gsr_raw);

% 2. Estrazione Componente Fasica (SCR)
GSR_fasica = gsr_raw - GSR_tonica; % Sottrazione

% In alternativa, si puo' filtrare direttamente
% in passa-banda la fasica
fc_scr = 1; % Hz (es. da 0.05 a 1 Hz)
[b_fas, a_fas] = butter(2, fc_scr/(fs/2), 'low');
GSR_fasica_filt = filtfilt(b_fas, a_fas, GSR_fasica);

% 3. Individuazione dei picchi emotivi
[peaks_eda, locs_eda] = islocalmax(GSR_fasica_filt,
'MinProminence', 0.01);
\end{lstlisting}

\section{Audio e Aliasing (Sottocampionamento)}
Per testare gli effetti del teorema di Nyquist-Shannon è utile procedere con la decimazione (downsampling) di un segnale audio pre-registrato.
\begin{lstlisting}[language=Matlab, frame=single, commentstyle=\color{green!50!black}, keywordstyle=\color{blue}]
% Caricamento audio originale
% (fs = 44100 Hz solitamente)
[y, fs_orig] = audioread('brano.wav');

% Fattore di sottocampionamento
M = 4;
y_down = y(1:M:end); % Salto M-1 campioni
fs_nuova = fs_orig / M; % Nuova frequenza di camp.

% Riproduzione: si notera' l'aliasing se
% fs_nuova < 2*f_max_segnale
sound(y_down, fs_nuova);
\end{lstlisting}

\section{Machine Learning Pratico (SVM)}
Nell'ecosistema MATLAB l'applicazione di una SVM linearmente o non linearmente separabile è agevolata da funzioni built-in.
\begin{lstlisting}[language=Matlab, frame=single, commentstyle=\color{green!50!black}, keywordstyle=\color{blue}]
% Addestramento di una Support Vector Machine
% X = matrice delle features (N samples x M features)
% Y = array delle labels (N samples)
svm_model = fitcsvm(X, Y, 'KernelFunction', 'linear',
'Standardize', true);

% Estrazione dei vettori di supporto
sv = svm_model.SupportVectors;

% Cross Validation (Holdout o K-fold)
cv_model = crossval(svm_model, 'KFold', 5);

% Stima dell'errore (o loss) generale
loss = kfoldLoss(cv_model);
fprintf('Accuratezza stimata: %.2f%%\n', (1-loss)*100);

% Predizione su nuovi dati
Y_pred = predict(svm_model, X_test);
\end{lstlisting}

